
\documentclass[11pt]{article}
\usepackage[margin=1in]{geometry}

% Core packages
\usepackage{amsmath,amssymb,amsthm}

% Paragraphs
\setlength{\parindent}{0pt}
\setlength{\parskip}{0.30\baselineskip}

\newtheorem{proposition}{Proposition}

\title{A Counterexample to the Jacobian Conjecture in Dimension Three}
\author{}
\date{\today}

\begin{document}
\maketitle

\section{The counterexample}

The Jacobian conjecture asserts that a polynomial map
$F\colon \mathbb{C}^n\to\mathbb{C}^n$ with constant nonzero Jacobian
determinant must be a polynomial automorphism.  The following map gives a
counterexample in dimension three.

\begin{proposition}
The polynomial map $F\colon\mathbb{C}^3\to\mathbb{C}^3$ defined by
\begin{align*}
F(x,y,z)=\bigl(&
 (1+xy)^3z+y^2(1+xy)(4+3xy),\\
&y+3x(1+xy)^2z+3xy^2(4+3xy),\\
&2x-3x^2y-x^3z
\bigr)
\end{align*}
has constant Jacobian determinant $-2$, but it is not injective.
Consequently, $F$ is not a polynomial automorphism.
\end{proposition}

\begin{proof}
Direct differentiation and simplification give
\[
\left|
\frac{\partial}{\partial(x,y,z)}
\begin{pmatrix}
(1+xy)^3z+y^2(1+xy)(4+3xy)\\
y+3x(1+xy)^2z+3xy^2(4+3xy)\\
2x-3x^2y-x^3z
\end{pmatrix}
\right|=-2.
\]
On the other hand, the three distinct points
\[
\left(0,0,-\frac14\right),\qquad
\left(1,-\frac32,\frac{13}{2}\right),\qquad
\left(-1,\frac32,\frac{13}{2}\right)
\]
have the same image, since
\begin{align*}
F\left(0,0,-\frac14\right)
&=\left((1+0)^3\left(-\frac14\right)+0,\;0+0+0,\;0-0-0\right)\\
&=\left(-\frac14,0,0\right),\\
F\left(1,-\frac32,\frac{13}{2}\right)
&=\left(-\frac{13}{16}+\frac9{16},\;
-\frac32+\frac{39}{8}-\frac{27}{8},\;
2+\frac92-\frac{13}{2}\right)\\
&=F\left(-1,\frac32,\frac{13}{2}\right)
=\left(-\frac14,0,0\right).
\end{align*}
Thus $F$ is not injective, despite having constant nonzero Jacobian
determinant.
\end{proof}

\paragraph{Acknowledgments.}
Thanks to my close friend Akhil for asking about the Jacobian conjecture and
to my close friend Fable for working on this during the World Cup final.

\end{document}